% Apunts de teoria en PDF · 1r ESO · Unitat 5: Decimals i proporcionalitat
% Mateix estil que els apunts de Batxillerat (franja i caixes de teoria en blau, exemples en roig), però més visuals,
% curts i amb dibuixos, per a l'alumnat de 3r ESO. Es compila amb pdflatex dues vegades (per als enllaços de l'índex),
% amb logo-fentmates.pdf (_eines/logo/pdf/horitzontal-clar.pdf) al costat.
\documentclass[11pt,a4paper]{article}
\usepackage[utf8]{inputenc}
\usepackage[T1]{fontenc}
\usepackage[catalan]{babel}
\usepackage[sfdefault,book,lining]{FiraSans}
\usepackage[italic,noplusnominus,noequal]{mathastext}
\linespread{1.18}
\usepackage[a4paper,top=2.4cm,bottom=2.2cm,left=1.6cm,right=1.6cm,headheight=1.1cm,headsep=0.5cm,footskip=1cm]{geometry}
\usepackage{amsmath,amssymb,cancel}
\usepackage{icomma}
\usepackage{graphicx}
\usepackage[table]{xcolor}
\usepackage{tikz}\usetikzlibrary{calc,positioning,arrows.meta,decorations.pathreplacing,shapes.geometric}
\usepackage{fontawesome5}
\usepackage{ccicons}
\usepackage{fancyhdr}
\usepackage{enumitem}
\usepackage{array}
\usepackage{needspace}
\usepackage[most]{tcolorbox}
\usepackage{titlesec}
\usepackage[hidelinks]{hyperref}
\hypersetup{pdftitle={Apunts de teoria. Unitat 5: Decimals i proporcionalitat. 1r ESO},pdfsubject={Matemàtiques 1r ESO},pdfkeywords={nombres decimals, dècimes, centèsimes, mil·lèsimes, comparar, arrodonir, truncar, operacions amb decimals, decimals i fraccions, raons, proporcions, proporcionalitat directa, proporcionalitat inversa, regla de tres, percentatges, augments i disminucions}}

% colors de la web (bloc d'Àlgebra: blau)
\definecolor{fmnavy}{HTML}{1B2748}
\definecolor{fmblue}{HTML}{3D6AA0}
\definecolor{fmbluel}{HTML}{4E81BD}
\definecolor{fmyellow}{HTML}{FFC94D}
\definecolor{fmgold}{HTML}{E9A919}
\definecolor{fmaccent}{HTML}{FF7A50}
\definecolor{fmdef}{HTML}{1D4F73}\definecolor{fmdefbg}{HTML}{EEF4F9}
\definecolor{fmex}{HTML}{A3203A}\definecolor{fmexbg}{HTML}{FBF0F2}
\definecolor{fmsf}{HTML}{C2410C}
\definecolor{fmok}{HTML}{15803D}
\definecolor{fmmuted}{HTML}{5D6685}
\definecolor{fmsurface}{HTML}{EEF1F8}
\definecolor{fmA}{HTML}{2F6FD6}\definecolor{fmB}{HTML}{EA6A10}\definecolor{fmR}{HTML}{E11D48}\definecolor{fmK}{HTML}{16A34A}\definecolor{fmV}{HTML}{7C3AED}

% lletres
\newcommand{\comic}{\fontfamily{ComicNeue-TLF}\selectfont}
\newcommand{\heavy}{\fontseries{eb}\selectfont}
\newcommand{\semi}{\fontseries{sb}\selectfont}
\DeclareUnicodeCharacter{2212}{-}

\setlength{\parindent}{0pt}
\setlength{\parskip}{7pt}
\setlength{\lineskip}{4pt}\setlength{\lineskiplimit}{4pt}
\setlength{\emergencystretch}{3em}
\setlist{itemsep=3pt,topsep=3pt,leftmargin=1.4em}
\setcounter{secnumdepth}{1}
\renewcommand{\thesection}{5.\arabic{section}}
\titleformat{\section}[hang]{\color{fmnavy}\fontsize{23}{27}\heavy\selectfont}{\tcbox[on line,colback=fmblue,colframe=fmblue,arc=3mm,boxsep=2pt,left=6pt,right=6pt,top=3pt,bottom=3pt]{\color{white}\thesection}}{0.45em}{}[\vspace{5pt}{\color{fmblue}\titlerule[2pt]}]
\titlespacing*{\section}{0pt}{4pt}{6pt}
\titleformat{\subsection}{\color{fmblue}\Large\semi}{}{0pt}{\faCaretRight\ }
\titlespacing*{\subsection}{0pt}{8pt}{2pt}

% frase de cada apartat, «a mà»
\newcommand{\frase}[1]{\par\vspace{-2pt}{\comic\Large\color{fmbluel}\faCommentDots[regular]\ \ #1\par}\vspace{2pt}}
\newcommand{\seccio}[3]{\Needspace{16\baselineskip}\section{#2}\hypertarget{#1}{}\frase{#3}}
% notes a mà dins dels exemples
\newcommand{\nota}[1]{{\comic\small\color{fmex!85!black}#1}}
\newcommand{\notab}[1]{{\comic\small\color{fmblue}#1}}
% bé / malament
\newcommand{\be}[1]{{\color{fmok}\faCheckCircle}\ \ensuremath{#1}}
\newcommand{\mal}[1]{{\color{fmR}\faTimesCircle}\ {\color{fmR}\ensuremath{#1}}}
% pas numerat
\newcommand{\pas}[1]{\tikz[baseline=(p.base)]\node[circle,fill=fmblue,text=white,inner sep=1.2pt,minimum size=1.25em,font=\small\heavy](p){#1};}
% etiqueta de color
\newcommand{\chip}[2]{\tcbox[on line,colback=#1!14,colframe=#1!60,arc=2mm,boxsep=1pt,left=5pt,right=5pt,top=1pt,bottom=1pt]{#2}}
\newcommand{\fmkey}[1]{\colorbox{fmyellow!55}{\strut #1}}
\newcommand{\cA}[1]{{\color{fmA}#1}}\newcommand{\cB}[1]{{\color{fmB}#1}}\newcommand{\cK}[1]{{\color{fmK}#1}}\newcommand{\cR}[1]{{\color{fmR}#1}}\newcommand{\cV}[1]{{\color{fmV}#1}}

% caixes
\tcbset{fmbase/.style={enhanced,arc=3mm,boxrule=1.2pt,left=4mm,right=4mm,top=2mm,bottom=2.5mm,before skip=8pt,after skip=8pt,fonttitle=\heavy\large,before upper=\raggedright}}
\newtcolorbox{idea}[1][Idea clau]{fmbase,colback=fmdefbg,colframe=fmdef,coltitle=white,title={\faKey\ \ #1}}
\newtcolorbox{exemple}[1][Exemple]{fmbase,colback=fmexbg,colframe=fmex,colbacktitle=fmexbg,coltitle=fmex,title={\faPencil*\ \ #1},titlerule=0.5pt,titlerule style={fmex!40,dashed}}
\newtcolorbox{compte}[1][Compte!]{fmbase,colback=fmsf!6,colframe=fmsf,coltitle=white,fonttitle=\comic\bfseries\Large,title={\faExclamationTriangle\ \ #1}}
\newtcolorbox{truc}[1][Truc]{fmbase,colback=fmyellow!20,colframe=fmgold,colbacktitle=fmyellow,coltitle=fmnavy,fonttitle=\comic\bfseries\Large,title={\faLightbulb[regular]\ \ #1}}

% dibuix de la targeta de la unitat, en blanc
\newcommand{\fmcard}{\begin{tikzpicture}[x=0.062cm,y=-0.062cm,white,line width=3pt,line cap=round]
  \begin{scope}[transparency group,opacity=0.15]
  \node[anchor=base west,inner sep=0pt,font=\fontsize{30}{30}\selectfont] at (14,44) {$2{,}35\quad 20\,\%$};
  \node[anchor=base west,inner sep=0pt,font=\fontsize{24}{24}\selectfont] at (14,100) {$\frac{3}{5}=\frac{12}{x}$};
  \end{scope}
\end{tikzpicture}}
\newcommand{\fmbrand}{\raisebox{-0.25\height}{\includegraphics[height=0.5cm]{logo-fentmates.pdf}}}

% regla de Ruffini: #1 = a, #2 = coeficients, #3 = fila del mig (sense el primer), #4 = fila de baix, #5 = nombre de coeficients
\newcommand{\ruffini}[5]{\begin{tikzpicture}[x=1.25cm,y=0.85cm,every node/.style={font=\large}]
  \foreach \v [count=\i] in {#2} \node (t\i) at (\i,0) {$\v$};
  \foreach \v [count=\i] in {#3} {\pgfmathtruncatemacro\j{\i+1}\node (m\j) at (\j,-1) {$\v$};}
  \foreach \v [count=\i] in {#4} \node (b\i) at (\i,-2.1) {$\v$};
  \node[text=fmB] at (0,-1) {$#1$};
  \draw[line width=1.1pt,fmnavy] (0.45,0.5)--(0.45,-1.5) (0.45,-1.5)--(#5+0.5,-1.5);
  \node[draw=fmR,line width=1.3pt,rounded corners=2mm,minimum size=0.75cm] at (#5,-2.1) {};
  \draw[-{Latex[length=2.2mm]},fmK,line width=1.1pt] (t1.south) -- (b1.north);
  \pgfmathtruncatemacro\last{#5-1}
  \foreach \i in {1,...,\last} {\pgfmathtruncatemacro\j{\i+1}
    \draw[-{Latex[length=2.2mm]},fmB,line width=1.1pt,shorten >=1pt] (b\i.north east) -- (m\j.south west);
    \draw[-{Latex[length=2.2mm]},fmA!70,line width=1pt,dashed,shorten >=2pt] ([xshift=8pt]t\j.south) -- ([xshift=8pt]b\j.north);}
\end{tikzpicture}}

% capçalera i peu
\pagestyle{fancy}\fancyhf{}
\fancyhead[L]{\small\fmbrand}
\fancyhead[R]{\small\color{fmmuted}1r ESO · Unitat 5: Decimals i proporcionalitat}
\fancyfoot[L]{\scriptsize\color{fmmuted}\href{https://creativecommons.org/licenses/by-nc-sa/4.0/deed.ca}{\ccbyncsa}\ \ FentMates · fentmates.org · CC BY-NC-SA 4.0}
\fancyfoot[R]{\small\color{fmmuted}\thepage}
\renewcommand{\headrulewidth}{0.4pt}\renewcommand{\headrule}{\color{fmmuted!40}\hrule height \headrulewidth}
\fancypagestyle{primera}{\fancyhf{}\renewcommand{\headrulewidth}{0pt}
  \fancyfoot[L]{\scriptsize\color{fmmuted}\href{https://creativecommons.org/licenses/by-nc-sa/4.0/deed.ca}{\ccbyncsa}\ \ FentMates · fentmates.org · CC BY-NC-SA 4.0}
  \fancyfoot[R]{\small\color{fmmuted}\thepage}}

% targeta de l'índex
\newcommand{\itile}[4]{\begin{tcolorbox}[enhanced,colback=white,colframe=fmblue!35,arc=3mm,boxrule=1pt,boxsep=0pt,left=6pt,right=6pt,top=5pt,bottom=5pt,before skip=0pt,after skip=7pt]
  \hyperlink{#1}{\begin{minipage}[c]{0.9cm}\centering{\color{fmblue}\LARGE #3}\end{minipage}\hspace{6pt}%
  \begin{minipage}[c]{0.78\linewidth}{\color{fmblue}\heavy #2}\\[-2pt]{\color{fmnavy}\semi\large #4}\end{minipage}}\end{tcolorbox}}

% taula amb vores arrodonides: el contingut es retalla amb la forma (sense @{} a les vores ni entre columnes)
\newsavebox\tabclipbox
\newcommand{\taulaclip}[3][3mm]{\sbox\tabclipbox{#3}\begin{tikzpicture}[baseline=(tcn.base)]
  \begin{scope}\clip[rounded corners=#1] (0,-\dp\tabclipbox) rectangle (\wd\tabclipbox,\ht\tabclipbox);
  \node[inner sep=0pt,outer sep=0pt,anchor=base west] (tcn) at (0,0) {\usebox\tabclipbox};\end{scope}
  \draw[#2,rounded corners=#1] (0,-\dp\tabclipbox) rectangle (\wd\tabclipbox,\ht\tabclipbox);
\end{tikzpicture}}
\newcommand{\taularod}[3][fmblue]{\taulaclip{#1,line width=1.2pt}{\renewcommand{\arraystretch}{1.5}\begin{tabular}{#2}#3\end{tabular}}}
% balança
\newcommand{\balanca}[3]{\begin{tikzpicture}[scale=#3,transform shape,every node/.style={anchor=center}]
  \fill[fmnavy] (-0.45,-1.9) rectangle (0.45,-1.75);
  \fill[fmnavy!80] (-0.12,-1.75) rectangle (0.12,0);
  \fill[fmgold] (0,0.05) circle (0.13);
  \draw[fmnavy,line width=2.4pt,line cap=round] (-2,0) -- (2,0);
  \foreach \s in {-1,1} {
    \draw[fmmuted,line width=0.8pt] (\s*2,0) -- (\s*2-0.75,-1.05) (\s*2,0) -- (\s*2+0.75,-1.05);
    \fill[fmyellow!70,draw=fmgold,line width=1pt] (\s*2-0.95,-1.05) arc[start angle=180,end angle=360,x radius=0.95,y radius=0.32] -- cycle;}
  \node[font=\Large,fill=fmA!15,draw=fmA,rounded corners=1.5mm,inner sep=3pt] at (-2,-0.68) {$#1$};
  \node[font=\Large,fill=fmB!15,draw=fmB,rounded corners=1.5mm,inner sep=3pt] at (2,-0.68) {$#2$};
\end{tikzpicture}}
% eixos amb quadrícula (dins d'un tikzpicture): xmin, xmax, ymin, ymax
\newcommand{\eixos}[4]{\draw[fmmuted!22,line width=0.4pt] (#1,#3) grid (#2,#4);
  \draw[-{Latex[length=2mm]},fmmuted,line width=0.8pt] (#1-0.2,0)--(#2+0.4,0) node[right,font=\small]{$x$};
  \draw[-{Latex[length=2mm]},fmmuted,line width=0.8pt] (0,#3-0.2)--(0,#4+0.4) node[above,font=\small]{$y$};}
% taula de traduccions (diccionari) amb les vores arrodonides
\newcommand{\dic}[2]{#1 & {\large$#2$}\\}
\newcommand{\dicratio}{0.62}
\newcommand{\dichead}{\begin{tabular}{@{}>{\raggedright}p{\dimexpr\dicratio\linewidth-12pt}>{\centering\arraybackslash}p{\dimexpr\linewidth-\dicratio\linewidth-12pt}@{}}\faComment[regular]\ Ho diem així\dots & \dots\ així\end{tabular}}
\newenvironment{dicbox}{\begin{tcolorbox}[enhanced,arc=3mm,boxrule=1.2pt,colframe=fmblue,colback=white,coltitle=white,fonttitle=\heavy,
  title={\dichead},left=0pt,right=0pt,top=0pt,bottom=1.5mm,boxsep=0pt,lefttitle=6pt,righttitle=6pt,before skip=0pt,after skip=0pt]\rowcolors{1}{white}{fmsurface}%
  \begin{tabular}{@{\hspace{6pt}}>{\raggedright}p{\dimexpr\dicratio\linewidth-12pt}>{\centering\arraybackslash}p{\dimexpr\linewidth-\dicratio\linewidth-12pt}@{\hspace{6pt}}}}
  {\end{tabular}\end{tcolorbox}}
% resum
\newcommand{\rtile}[3]{\begin{tcolorbox}[enhanced,colback=white,colframe=fmblue!45,arc=3mm,boxrule=1pt,left=3mm,right=3mm,top=2.5mm,bottom=2.5mm,fontupper=\large,
  title={\tcbox[on line,colback=fmblue,colframe=fmblue,arc=2mm,boxsep=0pt,left=4pt,right=4pt,top=1pt,bottom=1pt]{\color{white}\small\heavy #1}\ \ #2},coltitle=fmnavy,colbacktitle=fmdefbg,fonttitle=\semi,toptitle=0.5mm,bottomtitle=0.5mm]#3\end{tcolorbox}}
% portada: #1 subtítol a mà
\newcommand{\portada}[1]{\thispagestyle{primera}
\begin{tikzpicture}[remember picture,overlay]
  \fill[fmblue] (current page.north west) rectangle ([yshift=-9.4cm]current page.north east);
  \begin{scope}
    \clip (current page.north west) rectangle ([yshift=-9.4cm]current page.north east);
    \node[anchor=north west,inner sep=0pt] at ([xshift=6.6cm,yshift=-1.3cm]current page.north west) {\fmcard};
  \end{scope}
  \node[anchor=north west,inner sep=0pt] at ([xshift=1.6cm,yshift=-1.1cm]current page.north west) {\tcbox[colback=white,colframe=white,arc=2.5mm,boxsep=0pt,left=7pt,right=7pt,top=5pt,bottom=5pt]{\includegraphics[height=0.75cm]{logo-fentmates.pdf}}};
\end{tikzpicture}
\vspace*{0.55cm}
\tcbox[colback=white!15!fmblue,colframe=white!15!fmblue,arc=2.5mm,boxsep=0pt,left=6pt,right=6pt,top=3pt,bottom=3pt]{\color{white}\small\bfseries Apunts de teoria · 1r ESO · Matemàtiques}

\vspace{2pt}
{\raggedright\color{white}\fontsize{30}{36}\heavy\selectfont{\color{fmyellow}Unitat 5}\\[2pt]Decimals i proporcionalitat\par}
\vspace{8pt}
{\color{white}\comic\Large #1\par}
\vspace*{\dimexpr 7.4cm-\pagetotal\relax}
}
\newcommand{\llegenda}{\par\vspace{16pt}
{\color{fmnavy}\heavy\Large\faPalette\ \ Com es llegeixen aquests apunts}\par\smallskip
\begin{tcbraster}[raster columns=4,raster equal height,raster column skip=8pt,raster before skip=12pt]
\tcbox[colback=fmdefbg,colframe=fmdef,arc=3mm,boxrule=1.2pt,halign=center,width=\linewidth]{{\color{fmdef}\Large\faKey}\\[2pt]{\color{fmdef}\heavy Idea clau}\\[-1pt]{\small el que cal saber}}
\tcbox[colback=fmexbg,colframe=fmex,arc=3mm,boxrule=1.2pt,halign=center,width=\linewidth]{{\color{fmex}\Large\faPencil*}\\[2pt]{\color{fmex}\heavy Exemple}\\[-1pt]{\small pas a pas}}
\tcbox[colback=fmsf!6,colframe=fmsf,arc=3mm,boxrule=1.2pt,halign=center,width=\linewidth]{{\color{fmsf}\Large\faExclamationTriangle}\\[2pt]{\color{fmsf}\heavy Compte!}\\[-1pt]{\small errors típics}}
\tcbox[colback=fmyellow!20,colframe=fmgold,arc=3mm,boxrule=1.2pt,halign=center,width=\linewidth]{{\color{fmgold!80!black}\Large\faLightbulb[regular]}\\[2pt]{\color{fmnavy}\heavy Truc}\\[-1pt]{\small per a recordar-ho}}
\end{tcbraster}}
\newcommand{\resumtitol}{\clearpage
\hypertarget{resum}{}
{\color{fmnavy}\fontsize{23}{27}\heavy\selectfont\faClipboardList\ \ En resum}\par\vspace{5pt}{\color{fmblue}\titlerule[2pt]}
\frase{Tota la unitat en una pàgina.}}
% els 5 passos per a resoldre problemes: #2 = text del pas 2
\newcommand{\passosproblema}[2]{\begin{center}
\begin{tikzpicture}[s/.style={draw=fmblue,fill=fmdefbg,rounded corners=3mm,line width=1.2pt,minimum width=2.75cm,minimum height=2.45cm,align=center,font=\small},
  ic/.style={font=\LARGE,text=fmblue}]
  \foreach \i/\ic/\t in {1/\faBookOpen/{\textbf{Llegim}\\i entenem},2/\faQuestionCircle/{#1},3/\faEquals/{#2},4/\faCogs/{\textbf{Resolem}},5/\faCheckDouble/{\textbf{Comprovem}\\i responem}} {
    \node[s] (s\i) at (3.3*\i,0) {};
    \node[ic] at ([yshift=-0.5cm]s\i.north) {\ic};
    \node[align=center,font=\small,anchor=north] at ([yshift=-0.95cm]s\i.north) {\t};
    \node[circle,fill=fmblue,text=white,font=\small\heavy,inner sep=1.5pt,minimum size=1.3em] at (s\i.north west) {\i};}
  \foreach \i [evaluate=\i as \j using int(\i+1)] in {1,...,4} \draw[-{Latex[length=3mm]},fmmuted,line width=1.5pt] (s\i) -- (s\j);
\end{tikzpicture}
\end{center}}

\newcommand{\lo}{\mathopen{]}}\newcommand{\ro}{\mathclose{[}}
% punts ple i buit, nombres dels eixos i taula de valors
\newcommand{\pple}[2][fmblue]{\fill[#1] #2 circle (3.6pt);}
\newcommand{\nums}[4]{\foreach \k in {#1,...,#2} {\ifnum\k=0\else\node[below,font=\scriptsize,text=fmmuted] at (\k,0) {$\k$};\fi}
  \foreach \k in {#3,...,#4} {\ifnum\k=0\else\node[left,font=\scriptsize,text=fmmuted] at (0,\k) {$\k$};\fi}}
\newcommand{\tv}[3][fmblue]{\taulaclip{#1,line width=1.2pt}{\renewcommand{\arraystretch}{1.45}%
  \begin{tabular}{>{\columncolor{#1}\color{white}\heavy}c*{12}{w{c}{0.8cm}}}#2\\\rowcolor{#1!8}\cellcolor{#1}#3\end{tabular}}}
% etiqueta d'una recta o corba
\newcommand{\etq}[3]{\node[fill=white,inner sep=1.5pt,text=#1,font=\small] at #2 {#3};}


\def\intvdot#1#2#3{\ifx#2o\filldraw[fill=white,draw=#3,line width=1.3pt] (#1,0) circle (3.3pt);\fi\ifx#2c\fill[#3] (#1,0) circle (3.5pt);\fi}
\newcommand{\intv}[5]{\begin{tikzpicture}[x=0.4cm,baseline=-3pt]
  \draw[fmmuted,line width=0.8pt] (-6,0)--(7,0);
  \foreach \k in {-5,...,6} \draw[fmmuted,line width=0.6pt] (\k,-0.08)--(\k,0.08);
  \ifx#3a\def\xa{-6}\else\def\xa{#1}\fi
  \ifx#4a\def\xb{7}\else\def\xb{#2}\fi
  \ifx#4a\draw[#5,line width=2.6pt,-{Latex[length=2.6mm]}] (\xa,0)--(\xb,0);\else\ifx#3a\draw[#5,line width=2.6pt,{Latex[length=2.6mm]}-] (\xa,0)--(\xb,0);\else\draw[#5,line width=2.6pt] (\xa,0)--(\xb,0);\fi\fi
  \intvdot{#1}{#3}{#5}\intvdot{#2}{#4}{#5}
  \ifx#3a\else\node[below=2pt,font=\scriptsize] at (#1,0) {$#1$};\fi
  \ifx#4a\else\node[below=2pt,font=\scriptsize] at (#2,0) {$#2$};\fi
\end{tikzpicture}}
\newcommand{\passos}[1]{\begin{tcolorbox}[fmbase,colback=fmdefbg,colframe=fmdef,coltitle=white,title={\faListOl\ \ Com es resolen}]#1\end{tcolorbox}}
\newcommand{\passosT}[2]{\begin{tcolorbox}[fmbase,colback=fmdefbg,colframe=fmdef,coltitle=white,title={\faListOl\ \ #1}]#2\end{tcolorbox}}
\newcommand{\barra}[3]{\filldraw[fill=#3!35,draw=#3,line width=0.9pt] (#1-0.32,0) rectangle (#1+0.32,#2);}
\newcommand{\caixa}[7]{\draw[#7,line width=1.3pt] (#2,#1) -- (#3,#1) (#5,#1) -- (#6,#1) (#2,#1-0.25) -- (#2,#1+0.25) (#6,#1-0.25) -- (#6,#1+0.25);
  \filldraw[fill=#7!20,draw=#7,line width=1.3pt] (#3,#1-0.45) rectangle (#5,#1+0.45); \draw[#7,line width=2.2pt] (#4,#1-0.45) -- (#4,#1+0.45);}
\newcommand{\recta}[2]{\draw[-{Latex[length=2.5mm]},fmnavy,line width=1.1pt] (#1-0.6,0) -- (#2+0.6,0);
  \foreach \k in {#1,...,#2} {\draw[fmnavy,line width=0.9pt] (\k,-0.12) -- (\k,0.12); \node[below=3pt,font=\small] at (\k,0) {$\k$};}}
\newcommand{\pc}{\,\%}
\newcommand{\cent}[2]{\begin{tikzpicture}[x=#2,y=#2]
  \foreach \k in {0,...,99} {\pgfmathtruncatemacro\i{mod(\k,10)}\pgfmathtruncatemacro\j{9-int(\k/10)}
    \ifnum\k<#1 \filldraw[fill=fmB!45,draw=white,line width=0.6pt] (\i,\j) rectangle ++(1,1);
    \else \filldraw[fill=fmsurface,draw=white,line width=0.6pt] (\i,\j) rectangle ++(1,1);\fi}
  \draw[fmB,line width=1pt,rounded corners=1pt] (0,0) rectangle (10,10);
\end{tikzpicture}}


\begin{document}
\portada{Preus, mesures i repartiments: decimals, regles de tres i percentatges.}

{\color{fmnavy}\heavy\Large\faMapSigns\ \ Què aprendràs?}\par\smallskip
\begin{minipage}[t]{0.485\linewidth}
\itile{sec1}{5.1}{\faIcon{dot-circle}}{Nombres decimals}
\itile{sec2}{5.2}{\faSortAmountUp}{Comparar i aproximar}
\itile{sec3}{5.3}{\faPlus}{Sumar, restar i multiplicar}
\itile{sec4}{5.4}{\faDivide}{Dividir amb decimals}
\itile{sec5}{5.5}{\faIcon{exchange-alt}}{Decimals i fraccions}
\itile{sec6}{5.6}{\faBalanceScale}{Raons i proporcions}
\end{minipage}\hfill
\begin{minipage}[t]{0.485\linewidth}
\itile{sec7}{5.7}{\faArrowUp}{Proporcionalitat directa}
\itile{sec8}{5.8}{\faArrowDown}{Proporcionalitat inversa}
\itile{sec9}{5.9}{\faPercent}{Percentatges}
\itile{sec10}{5.10}{\faChartLine}{Augments i disminucions}
\itile{resum}{\phantom{5.10}}{\faClipboardList}{En resum}
\end{minipage}
\llegenda

% ============================================================ 5.1
\clearpage
\seccio{sec1}{Nombres decimals}{Quantitats que no són senceres: una part entera i una part decimal.}

\vfill
\begin{minipage}[c]{0.5\linewidth}
\begin{idea}[Què és un nombre decimal?]
Un preu de 2,35 €, una altura d'1,68 m, 36,6 °C\dots\\[4pt]
{\centering\fontsize{26}{30}\selectfont\cA{245},\cB{378}\par}\smallskip
\cA{\textbf{Part entera}} (245), a l'esquerra de la coma. \ \cB{\textbf{Part decimal}} (378): un tros d'unitat, menys que 1.
\end{idea}
\end{minipage}\hfill
\begin{minipage}[c]{0.47\linewidth}\centering
\begin{tikzpicture}[scale=0.3]
  \foreach \i in {0,...,9} \foreach \j in {0,...,9} {\pgfmathtruncatemacro\n{\i*10+(9-\j)}\ifnum\n<27 \fill[fmB!45] (\i,\j) rectangle ++(1,1);\fi}
  \draw[fmmuted,step=1] (0,0) grid (10,10); \draw[fmnavy,line width=1.2pt] (0,0) rectangle (10,10);
\end{tikzpicture}\\
{\comic\small\color{fmmuted}$0{,}27$: 2 columnes (dècimes) i 7 quadrets (centèsimes)}
\end{minipage}

\vfill
\begin{tcbraster}[raster columns=3,raster equal height,raster column skip=8pt]
\begin{tcolorbox}[fmbase,colback=fmA!8,colframe=fmA,coltitle=white,title={Dècima},before upper=\centering]
La unitat en 10 parts\par\smallskip {\large$0{,}1=\frac1{10}$}
\end{tcolorbox}
\begin{tcolorbox}[fmbase,colback=fmB!8,colframe=fmB,coltitle=white,title={Centèsima},before upper=\centering]
La unitat en 100 parts\par\smallskip {\large$0{,}01=\frac1{100}$}
\end{tcolorbox}
\begin{tcolorbox}[fmbase,colback=fmK!8,colframe=fmK,coltitle=white,title={Mil·lèsima},before upper=\centering]
La unitat en 1000 parts\par\smallskip {\large$0{,}001=\frac1{1000}$}
\end{tcolorbox}
\end{tcbraster}

\vfill
\begin{center}
\taularod{*{6}{>{\hspace{8pt}}c<{\hspace{8pt}}}}{
\rowcolor{fmblue}\color{white}\heavy C & \color{white}\heavy D & \color{white}\heavy U & \color{white}\heavy d & \color{white}\heavy c & \color{white}\heavy m\\
\cA{\Large 2} & \cA{\Large 4} & \cA{\Large 5}\rlap{\hspace{6pt}\Large\heavy ,} & \cB{\Large 3} & \cB{\Large 7} & \cB{\Large 8}\\
\rowcolor{fmsurface}\small 200 & \small 40 & \small 5 & \small 0,3 & \small 0,07 & \small 0,008
}
\end{center}

\vfill
\begin{tcbraster}[raster columns=2,raster equal height,raster column skip=10pt]
\begin{tcolorbox}[fmbase,colback=fmdefbg,colframe=fmdef,coltitle=white,title={Lectura}]
La part entera, «unitats», i la decimal amb el nom de l'última xifra: \ 245,378, «245 unitats i 378 mil·lèsimes».
\end{tcolorbox}
\begin{tcolorbox}[fmbase,colback=fmdefbg,colframe=fmdef,coltitle=white,title={Descomposició}]
Sumem el valor de cada xifra:\par\smallskip
$245{,}378=200+40+5+0{,}3+0{,}07+0{,}008$
\end{tcolorbox}
\end{tcbraster}

\vfill
\begin{compte}[Zeros a la dreta, sí; al mig, no]
$2{,}5=2{,}50=2{,}500$ \ (5 dècimes $=$ 50 centèsimes). \ Però $2{,}05\neq2{,}5$: el 5 ja no són dècimes, sinó centèsimes.
\end{compte}

\vfill
\begin{exemple}[Exemples]
$3{,}07$: «3 unitats i 7 centèsimes» \qquad $0{,}125$: «125 mil·lèsimes» \qquad 6 centèsimes $=0{,}06$\\[4pt]
$12{,}4=10+2+0{,}4$ \qquad $7{,}25=7+0{,}2+0{,}05$ \qquad 15 dècimes $=1{,}5$
\end{exemple}

\vfill
% ============================================================ 5.2
\clearpage
\seccio{sec2}{Comparar i aproximar decimals}{Primer la part entera; després, xifra a xifra.}

\vfill
\begin{minipage}[c]{0.52\linewidth}
\passosT{Com es comparen}{\pas{1} La \textbf{part entera}: és més gran el que la té més gran. \ $12{,}3>9{,}87$\\[3pt]
\pas{2} Si són iguals, les \textbf{dècimes}; després, les centèsimes\dots\\[3pt]
\pas{3} Truc: \textbf{igualem decimals} amb zeros: \ $3{,}5\ ?\ 3{,}48\ \to\ 3{,}50>3{,}48$}
\end{minipage}\hfill
\begin{minipage}[c]{0.45\linewidth}
\begin{compte}[Error típic]
\mal{3{,}48>3{,}5} perquè 48 és més que 5? \ No: $3{,}5$ té 5 dècimes i $3{,}48$ només 4.
\end{compte}
\end{minipage}

\vfill
\begin{center}
\begin{tikzpicture}[x=0.6cm]
  \draw[fmnavy,line width=1.1pt] (0,0) -- (10,0);
  \foreach \k in {0,...,10} \draw[fmnavy] (\k,-0.1) -- (\k,0.1);
  \node[below,font=\small] at (0,-0.1) {2}; \node[below,font=\small] at (10,-0.1) {3};
  \node[below,font=\scriptsize] at (7,-0.1) {2,7}; \node[below,font=\scriptsize] at (8,-0.1) {2,8};
  \fill[fmR] (7.4,0) circle (2.5pt);
  \draw[fmmuted,dashed] (7,-0.5) -- (1,-1.5); \draw[fmmuted,dashed] (8,-0.5) -- (11,-1.5);
  \draw[fmnavy,line width=1.1pt] (1,-1.6) -- (11,-1.6);
  \foreach \k in {1,...,11} \draw[fmnavy] (\k,-1.7) -- (\k,-1.5);
  \node[below,font=\scriptsize] at (1,-1.7) {2,70}; \node[below,font=\scriptsize] at (11,-1.7) {2,80};
  \fill[fmR] (5,-1.6) circle (2.5pt) node[below=3pt,font=\small\heavy,text=fmR] {2,74};
\end{tikzpicture}\\
{\comic\small\color{fmmuted}entre 2 i 3, de dècima en dècima; i fem zoom entre 2,7 i 2,8}
\end{center}

\vfill
\begin{truc}[Sempre n'hi ha més]
Entre dos decimals sempre hi ha infinits decimals: entre $1{,}4$ i $1{,}5$ hi ha $1{,}41$, $1{,}45$, $1{,}499$\dots
\end{truc}

\vfill
\begin{minipage}[c]{0.5\linewidth}
\begin{idea}[Aproximar]
\textbf{Truncar}: tallem les xifres de la dreta.\\[3pt]
\textbf{Arrodonir}: mirem la primera xifra que eliminem. 0, 1, 2, 3 o 4: l'última es queda igual; 5, 6, 7, 8 o 9: li sumem 1.
\end{idea}
\end{minipage}\hfill
\begin{minipage}[c]{0.47\linewidth}\centering\small
\taularod{>{\heavy}l*{2}{>{\hspace{6pt}}c<{\hspace{6pt}}}}{
\rowcolor{fmblue}\color{white}3,4682 a\dots & \color{white}\heavy Truncar & \color{white}\heavy Arrodonir\\
unitats & 3 & 3\\
\rowcolor{fmsurface}dècimes & 3,4 & \cR{3,5}\\
centèsimes & 3,46 & \cR{3,47}\\
\rowcolor{fmsurface}mil·lèsimes & 3,468 & 3,468
}
\end{minipage}

\vfill
\begin{exemple}[Exemples]
\begin{minipage}[t]{0.48\linewidth}
De menor a major: \ $0{,}9<1{,}09<1{,}9<1{,}91$\\[4pt]
$4{,}6\ ?\ 4{,}59$: \ $4{,}60>4{,}59$
\end{minipage}\hfill
\begin{minipage}[t]{0.49\linewidth}
Arrodoneix 7,856 a les dècimes: \ $\boxed{7{,}9}$\\[4pt]
Un preu de 3,4567 € en cèntims: \ $\boxed{3{,}46\text{ €}}$
\end{minipage}
\end{exemple}

\vfill
% ============================================================ 5.3
\clearpage
\seccio{sec3}{Sumar, restar i multiplicar decimals}{Comes alineades per a sumar; comptem decimals per a multiplicar.}

\vfill
\begin{minipage}[c]{0.5\linewidth}
\passosT{Sumar i restar}{\pas{1} Un davall de l'altre, amb les \textbf{comes alineades}.\\[3pt]
\pas{2} Completem amb \textbf{zeros} si cal.\\[3pt]
\pas{3} Operem com amb naturals i baixem la coma.}
\end{minipage}\hfill
\begin{minipage}[c]{0.22\linewidth}\centering\Large
$\begin{array}{r}12{,}50\\+\ \ 3{,}78\\\hline16{,}28\end{array}$
\end{minipage}\hfill
\begin{minipage}[c]{0.22\linewidth}\centering\Large
$\begin{array}{r}15{,}00\\-\ \ 6{,}47\\\hline8{,}53\end{array}$
\end{minipage}

\vfill
\begin{truc}[La coma amagada]
Un natural té la coma al final: \ $15=15{,}00$.
\end{truc}

\vfill
\begin{minipage}[c]{0.52\linewidth}
\begin{idea}[Multiplicar decimals]
Multipliquem sense comes i separem tantes xifres decimals com \textbf{sumen} les dels dos factors:\\[4pt]
$3{,}25\cdot1{,}4$: \ $325\cdot14=4550$\\
$2+1=3$ decimals: \ \boxed{4{,}550}
\end{idea}
\end{minipage}\hfill
\begin{minipage}[c]{0.45\linewidth}
\begin{compte}[Estima abans]
$3{,}25\cdot1{,}4$ és una mica més que $3\cdot1=3$ i menys que $4\cdot2=8$. Si et dona 45,5 o 0,455, la coma està malament.
\end{compte}
\end{minipage}

\vfill
\begin{tcbraster}[raster columns=2,raster equal height,raster column skip=10pt]
\begin{tcolorbox}[fmbase,colback=fmK!8,colframe=fmK,coltitle=white,title={$\times$ 10, 100, 1000\dots}]
La coma \textbf{a la dreta} tants llocs com zeros:\par\smallskip
$3{,}7\cdot100=370$ \qquad $0{,}052\cdot1000=52$
\end{tcolorbox}
\begin{tcolorbox}[fmbase,colback=fmB!8,colframe=fmB,coltitle=white,title={$:$ 10, 100, 1000\dots}]
La coma \textbf{a l'esquerra}:\par\smallskip
$45{,}2:1000=0{,}0452$ \qquad $8:10=0{,}8$
\end{tcolorbox}
\end{tcbraster}

\vfill
\begin{exemple}[Exemples]
\begin{minipage}[t]{0.48\linewidth}
$4{,}7+12{,}35+0{,}8=\boxed{17{,}85}$\\[4pt]
$20-7{,}63=\boxed{12{,}37}$\\[4pt]
$0{,}6\cdot0{,}3=\boxed{0{,}18}$ \ \nota{$6\cdot3=18$, 2 decimals}
\end{minipage}\hfill
\begin{minipage}[t]{0.49\linewidth}
$2{,}5\cdot4{,}2=\boxed{10{,}50}$\\[4pt]
$250\cdot0{,}1=25$ \ \nota{$\cdot0{,}1$ és $:10$}\\[4pt]
3 kg de pomes a 1,85 €: \ $3\cdot1{,}85=\boxed{5{,}55\text{ €}}$
\end{minipage}
\end{exemple}

\vfill
\begin{exemple}[Exemple: \ comprem 2,5 kg de taronges a 1,40 €/kg i una barra de pa d'1,15 €. Paguem amb 10 €. Quant ens tornen?]
Taronges: \ $2{,}5\cdot1{,}40=3{,}50$ € \qquad Total: \ $3{,}50+1{,}15=4{,}65$ € \qquad Canvi: \ $10-4{,}65=\boxed{5{,}35\text{ €}}$
\end{exemple}

\vfill
% ============================================================ 5.4
\clearpage
\seccio{sec4}{Dividir amb decimals}{Si el divisor té comes, les fem desaparéixer.}

\vfill
\begin{tcbraster}[raster columns=3,raster equal height,raster column skip=8pt]
\begin{tcolorbox}[fmbase,colback=fmA!8,colframe=fmA,coltitle=white,title={Decimal entre natural}]
Dividim com sempre; en baixar la primera xifra decimal, posem la \textbf{coma} al quocient.\par\smallskip
{\centering$9{,}75:3=3{,}25$\par}
\end{tcolorbox}
\begin{tcolorbox}[fmbase,colback=fmB!8,colframe=fmB,coltitle=white,title={Traure decimals}]
Si no és exacta, coma al quocient, un 0 al residu i continuem.\par\smallskip
{\centering$7:4=1{,}75$\par}
\end{tcolorbox}
\begin{tcolorbox}[fmbase,colback=fmK!8,colframe=fmK,coltitle=white,title={Entre un decimal}]
Multipliquem dividend i divisor per 10, 100\dots\ fins que el divisor siga natural.\par\smallskip
{\centering$12{,}6:0{,}35=1260:35$\par}
\end{tcolorbox}
\end{tcbraster}

\vfill
\begin{minipage}[c]{0.4\linewidth}\centering
\begin{tikzpicture}[font=\Large]
  \node at (0,0) {$7$}; \node at (1.3,0) {$4$};
  \draw[fmnavy,line width=1.2pt] (0.65,0.35) -- (0.65,-0.35) -- (2.2,-0.35);
  \node[text=fmA] at (1.4,-0.75) {$1{,}75$};
  \node at (0,-0.75) {$30$}; \node at (0,-1.45) {$20$}; \node[text=fmR] at (0,-2.15) {$0$};
\end{tikzpicture}
\end{minipage}\hfill
\begin{minipage}[c]{0.57\linewidth}
\begin{exemple}[Exemple: \ $7:4$ amb decimals]
$7:4=1$, residu 3. \ Coma i un 0: $30:4=7$, residu 2.\\[3pt]
Un altre 0: $20:4=5$, residu 0. \qquad \boxed{7:4=1{,}75}
\end{exemple}
\end{minipage}

\vfill
\begin{exemple}[Exemple: \ $12{,}6:0{,}35$]
El divisor té 2 decimals: multipliquem tots dos per 100: \ $1260:35=\boxed{36}$ \qquad \nota{el quocient no canvia}
\end{exemple}

\vfill
\begin{truc}[Dividir entre menys d'1 fa créixer]
$12{,}6:0{,}35=36$: en 12,6 hi caben moltes vegades 0,35. Té sentit!
\end{truc}

\vfill
\begin{exemple}[Exemples]
\begin{minipage}[t]{0.48\linewidth}
$15{,}6:4=\boxed{3{,}9}$\\[4pt]
$5:8=\boxed{0{,}625}$\\[4pt]
$4{,}5:0{,}9=45:9=\boxed{5}$
\end{minipage}\hfill
\begin{minipage}[t]{0.49\linewidth}
$2{,}4:0{,}06=240:6=\boxed{40}$\\[4pt]
$10:3=3{,}333\ldots\approx\boxed{3{,}33}$ \ \nota{no acaba mai}\\[4pt]
$7{,}2:1{,}2=72:12=\boxed{6}$
\end{minipage}
\end{exemple}

\vfill
\begin{exemple}[Exemple: \ 4 amics paguen 25,40 € per un sopar. I un pot de 1,5 L de suc s'omple amb gots de 0,25 L]
Cadascú: \ $25{,}40:4=\boxed{6{,}35\text{ €}}$ \qquad Gots: \ $1{,}5:0{,}25=150:25=\boxed{6\text{ gots}}$
\end{exemple}

\vfill
% ============================================================ 5.5
\clearpage
\seccio{sec5}{Decimals i fraccions}{Dividim el numerador entre el denominador.}

\vfill
\begin{tcbraster}[raster columns=3,raster equal height,raster column skip=8pt]
\begin{tcolorbox}[fmbase,colback=fmK!8,colframe=fmK,coltitle=white,title={Exacte},before upper=\centering]
Xifres decimals \textbf{finites}.\par\smallskip
{\Large$\dfrac38=0{,}375$}
\end{tcolorbox}
\begin{tcolorbox}[fmbase,colback=fmA!8,colframe=fmA,coltitle=white,title={Periòdic pur},before upper=\centering]
Es repeteixen \textbf{des de la coma}.\par\smallskip
{\Large$\dfrac{5}{11}=0{,}\overline{45}$}\par{\small$0{,}454545\dots$}
\end{tcolorbox}
\begin{tcolorbox}[fmbase,colback=fmV!8,colframe=fmV,coltitle=white,title={Periòdic mixt},before upper=\centering]
Abans del període, xifres que \textbf{no} es repeteixen.\par\smallskip
{\Large$\dfrac{7}{12}=0{,}58\overline{3}$}\par{\small$0{,}58333\dots$}
\end{tcolorbox}
\end{tcbraster}

\vfill
\begin{center}
\begin{tikzpicture}
  \node[font=\fontsize{36}{36}\selectfont] (n) {$\cB{2}{,}\cK{58}\overline{\cA{3}}$};
  \node[anchor=east,font=\large\heavy,text=fmB] (pe) at ([xshift=-1.6cm]n.west) {part entera};
  \node[anchor=north,font=\large\heavy,text=fmK!80!black] (ap) at ([xshift=-2pt,yshift=-0.9cm]n.south) {anteperíode};
  \node[anchor=west,font=\large\heavy,text=fmA] (p) at ([xshift=1.6cm]n.east) {període};
  \draw[-{Latex[length=3mm]},fmB,line width=1.3pt] (pe.east) -- ([xshift=4pt,yshift=-2pt]n.west);
  \draw[-{Latex[length=3mm]},fmK,line width=1.3pt] (ap.north) -- ([xshift=-2pt,yshift=4pt]n.south);
  \draw[-{Latex[length=3mm]},fmA,line width=1.3pt] (p.west) -- ([xshift=-2pt,yshift=2pt]n.east);
\end{tikzpicture}
\end{center}

\vfill
\begin{idea}
Els tres tipus són \textbf{racionals}. Els decimals amb infinites xifres \textbf{no periòdiques} (com $\pi=3{,}14159\dots$) són \textbf{irracionals}.
\end{idea}

\vfill
\subsection{Com saber-ho sense dividir}
{\large Simplifica la fracció i descompon el \textbf{denominador}:}

\vfill
\begin{center}
\begin{tikzpicture}[b/.style={draw=#1,fill=#1!10,line width=1.2pt,rounded corners=3mm,minimum width=5cm,minimum height=1.7cm,align=center}]
  \node[b=fmK] at (0,0) {Només \textbf{2} i/o \textbf{5}\\[2pt]$\dfrac38$: \ $8=2^3$};
  \node[b=fmA] at (5.6,0) {Ni 2 ni 5\\[2pt]$\dfrac{5}{11}$: \ $11$};
  \node[b=fmV] at (11.2,0) {2 o 5 \textbf{i} altres\\[2pt]$\dfrac{7}{12}$: \ $12=2^2\cdot3$};
  \node[font=\large\heavy,text=fmK!80!black] at (0,-1.25) {$\to$ exacte};
  \node[font=\large\heavy,text=fmA] at (5.6,-1.25) {$\to$ periòdic pur};
  \node[font=\large\heavy,text=fmV] at (11.2,-1.25) {$\to$ periòdic mixt};
\end{tikzpicture}
\end{center}

\vfill
\begin{minipage}[t]{0.5\linewidth}
\begin{truc}[Per què 2 i 5?]
Perquè $10=2\cdot5$: només les fraccions que es poden escriure amb denominador 10, 100, 1000\dots\ donen decimals exactes.
\end{truc}
\end{minipage}\hfill
\begin{minipage}[t]{0.47\linewidth}
\begin{compte}[Primer, simplifica!]
$\dfrac{21}{28}$: \ $28=2^2\cdot7$ sembla mixt\dots\\[3pt]
però $\dfrac{21}{28}=\dfrac34$ \ $\to$ \ \textbf{exacte} ($0{,}75$).
\end{compte}
\end{minipage}

\vfill
% ============================================================ 5.6
\clearpage
\seccio{sec6}{Raons i proporcions}{Comparem dues quantitats dividint-les.}

\vfill
\begin{minipage}[c]{0.52\linewidth}
\begin{idea}[Raó]
El quocient de dos nombres, $\dfrac ab$ («$a$ és a $b$»). Serveix per a \textbf{comparar} quantitats.
\end{idea}
\end{minipage}\hfill
\begin{minipage}[c]{0.45\linewidth}
\begin{exemple}[Exemple: \ 12 xics i 15 xiques]
$\dfrac{12}{15}=\dfrac45=0{,}8$\\[3pt]
Per cada 4 xics hi ha 5 xiques.
\end{exemple}
\end{minipage}

\vfill
\begin{truc}[Raó o fracció?]
S'assemblen, però els termes d'una raó poden ser decimals i comparen dues quantitats \textbf{diferents} (xics i xiques, km i hores\dots), no parts d'una unitat.
\end{truc}

\vfill
\begin{minipage}[c]{0.5\linewidth}
\begin{idea}[Proporció]
Una igualtat entre dues raons: «$a$ és a $b$ com $c$ és a $d$».\\[3pt]
{\centering\Large$\dfrac{\cB{a}}{\cA{b}}=\dfrac{\cA{c}}{\cB{d}}$\par}\smallskip
\cB{$a$ i $d$}: els \textbf{extrems}; \ \cA{$b$ i $c$}: els \textbf{mitjans}.
\end{idea}
\end{minipage}\hfill
\begin{minipage}[c]{0.47\linewidth}
\begin{tcolorbox}[enhanced,colback=fmyellow!18,colframe=fmgold,arc=4mm,boxrule=1.6pt,before upper=\centering,top=2mm,bottom=2mm]
\textbf{Propietat fonamental}\par\smallskip
{\large$a\cdot d=b\cdot c$}\par\smallskip
\small extrems per extrems igual a mitjans per mitjans (els productes creuats)
\end{tcolorbox}
\end{minipage}

\vfill
\begin{exemple}[Exemple: \ el quart proporcional]
$\dfrac35=\dfrac{12}{x}\ \Rightarrow\ 3\cdot x=5\cdot12\ \Rightarrow\ x=\dfrac{5\cdot12}{3}=\boxed{20}$ \qquad Comprovem: \ $3\cdot20=60=5\cdot12$ \be{}
\end{exemple}

\vfill
\begin{exemple}[Exemples]
\begin{minipage}[t]{0.48\linewidth}
És una proporció $\frac{4}{6}=\frac{10}{15}$? \ $60=60$ \be{}\\[4pt]
$\frac{x}{8}=\frac{15}{24}$: \ $24x=120$, \ \boxed{x=5}
\end{minipage}\hfill
\begin{minipage}[t]{0.49\linewidth}
Una recepta: 3 ous per 250 g de farina. Amb 1000 g? \ $\frac{3}{250}=\frac{x}{1000}$, \ \boxed{12\text{ ous}}\\[4pt]
Raó km/h: 180 km en 2 h, \ 90 km/h
\end{minipage}
\end{exemple}

\vfill
\begin{compte}[L'ordre de la raó]
La raó entre xics i xiques és $\frac{12}{15}$; entre xiques i xics, $\frac{15}{12}$. No són la mateixa!
\end{compte}

\vfill
% ============================================================ 5.7
\clearpage
\seccio{sec7}{Proporcionalitat directa}{El doble, el doble; la meitat, la meitat.}

\vfill
\begin{minipage}[c]{0.52\linewidth}
\begin{idea}[Magnituds directament proporcionals]
Si multipliquem (o dividim) una per un nombre, l'altra queda multiplicada (o dividida) pel \textbf{mateix} nombre.\\[3pt]
El \textbf{quocient} es manté: és la \textbf{constant de proporcionalitat}.
\end{idea}
\end{minipage}\hfill
\begin{minipage}[c]{0.45\linewidth}\centering
\tv{kg&1&2&3&5&10}{€&1,80&3,60&5,40&9&18}\\[4pt]
{\small$\dfrac{3{,}60}{2}=\dfrac{5{,}40}{3}=\dfrac95=1{,}80$ €/kg}
\end{minipage}

\vfill
\begin{compte}[No n'hi ha prou que «augmenten juntes»]
L'edat i l'altura augmenten juntes, però als 20 anys no fem el doble que als 10. I un taxi amb baixada de bandera: 10 km no costen el doble que 5.
\end{compte}

\vfill
\begin{exemple}[Exemple: \ 3 kg de taronges costen 5,40 €. Quant costen 7 kg?]
\begin{minipage}[t]{0.48\linewidth}
\textbf{\cA{Reducció a la unitat}}:\\[3pt]
1 kg: \ $5{,}40:3=1{,}80$ €\\[2pt]
7 kg: \ $1{,}80\cdot7=\boxed{12{,}60\text{ €}}$
\end{minipage}\hfill
\begin{minipage}[t]{0.49\linewidth}
\textbf{\cB{Regla de tres directa}}:\\[3pt]
$\begin{array}{rcl}3\text{ kg}&\to&5{,}40\text{ €}\\7\text{ kg}&\to&x\end{array}$ \quad $x=\dfrac{7\cdot5{,}40}{3}=\boxed{12{,}60\text{ €}}$
\end{minipage}
\end{exemple}

\vfill
\begin{truc}[Regla de tres directa]
Multipliquem \textbf{en creu} i dividim entre el que queda.
\end{truc}

\vfill
\begin{exemple}[Exemples]
\begin{minipage}[t]{0.48\linewidth}
Un cotxe gasta 6 L cada 100 km. En 350 km?\\ $\dfrac{350\cdot6}{100}=\boxed{21\text{ L}}$\\[4pt]
5 bolígrafs costen 3,50 €. I 8? \ $\boxed{5{,}60\text{ €}}$
\end{minipage}\hfill
\begin{minipage}[t]{0.49\linewidth}
Una impressora fa 12 còpies per minut. En 15 min?\\ $\boxed{180\text{ còpies}}$\\[4pt]
En un mapa, 2 cm són 50 km. I 7 cm? \ $\boxed{175\text{ km}}$
\end{minipage}
\end{exemple}

\vfill
\begin{exemple}[Exemple: \ una recepta per a 4 persones porta 300 g d'arròs i 600 mL d'aigua. I per a 6 persones?]
Per a 1 persona: \ 75 g i 150 mL \qquad Per a 6: \ $75\cdot6=\boxed{450\text{ g}}$ d'arròs i $150\cdot6=\boxed{900\text{ mL}}$ d'aigua
\end{exemple}

\vfill
% ============================================================ 5.8
\clearpage
\seccio{sec8}{Proporcionalitat inversa}{El doble, la meitat.}

\vfill
\begin{minipage}[c]{0.52\linewidth}
\begin{idea}[Magnituds inversament proporcionals]
Si multipliquem una per un nombre, l'altra queda \textbf{dividida} pel mateix nombre (i al revés).\\[3pt]
El \textbf{producte} es manté: és la constant de proporcionalitat inversa.
\end{idea}
\end{minipage}\hfill
\begin{minipage}[c]{0.45\linewidth}\centering
{\small Un viatge de 240 km:}\\[2pt]
\tv{km/h&40&60&80&120}{h&6&4&3&2}\\[4pt]
{\small$40\cdot6=60\cdot4=80\cdot3=240$}
\end{minipage}

\vfill
\begin{tcbraster}[raster columns=2,raster equal height,raster column skip=10pt]
\begin{tcolorbox}[fmbase,colback=fmK!8,colframe=fmK,coltitle=white,title={\faArrowUp\ \ Directa}]
Una puja, l'altra puja \textbf{igual}.\par\smallskip
Es manté el \textbf{quocient}. \ Regla: en creu.
\end{tcolorbox}
\begin{tcolorbox}[fmbase,colback=fmB!8,colframe=fmB,coltitle=white,title={\faArrowDown\ \ Inversa}]
Una puja, l'altra \textbf{baixa}: el doble, la meitat.\par\smallskip
Es manté el \textbf{producte}. \ Regla: en línia.
\end{tcolorbox}
\end{tcbraster}

\vfill
\begin{exemple}[Exemple: \ 4 pintors pinten una casa en 6 dies. Quant tardarien 3 pintors?]
\begin{minipage}[t]{0.48\linewidth}
\textbf{\cA{Reducció a la unitat}}:\\[3pt]
1 pintor tardaria 4 vegades més: \ $6\cdot4=24$ dies\\[2pt]
3 pintors: \ $24:3=\boxed{8\text{ dies}}$
\end{minipage}\hfill
\begin{minipage}[t]{0.49\linewidth}
\textbf{\cB{Regla de tres inversa}}: el producte es manté:\\[3pt]
$4\cdot6=3\cdot x$ \quad $x=\dfrac{24}{3}=\boxed{8\text{ dies}}$
\end{minipage}
\end{exemple}

\vfill
\begin{compte}[Directa, inversa o cap?]
Pregunta't: si una es fa el doble, què li passa a l'altra? \ El doble: \textbf{directa}. \ La meitat: \textbf{inversa}. \ Ni una cosa ni l'altra: \textbf{cap} (com l'edat i l'altura).
\end{compte}

\vfill
\begin{exemple}[Exemples]
\begin{minipage}[t]{0.48\linewidth}
3 aixetes omplin un dipòsit en 8 h. I 6 aixetes? \ $\frac{3\cdot8}{6}=\boxed{4\text{ h}}$\\[4pt]
A 90 km/h tardem 2 h. A 60 km/h? \ $\frac{90\cdot2}{60}=\boxed{3\text{ h}}$
\end{minipage}\hfill
\begin{minipage}[t]{0.49\linewidth}
Un pastís per a 8 persones dona trossos de 120 g. Per a 12? \ $\frac{8\cdot120}{12}=\boxed{80\text{ g}}$\\[4pt]
El preu del pa i els kg que compres: \textbf{directa}. \ Els obrers i els dies: \textbf{inversa}.
\end{minipage}
\end{exemple}

\vfill
\begin{truc}[Comprova amb el sentit comú]
Més pintors, menys dies; més velocitat, menys temps. Si el resultat va al revés, has fet una regla directa en lloc d'inversa.
\end{truc}

\vfill
% ============================================================ 5.9
\clearpage
\seccio{sec9}{Percentatges}{Un percentatge diu quantes parts prenem de cada 100.}

\vfill
\begin{minipage}[c]{0.3\linewidth}\centering
\cent{35}{0.42cm}\\[4pt]
{\comic\color{fmB}35 de cada 100}
\end{minipage}\hfill
\begin{minipage}[c]{0.67\linewidth}
\begin{idea}[Tres maneres de dir el mateix]
{\centering\fontsize{20}{26}\selectfont$\cB{35\pc}=\cA{\dfrac{35}{100}}=\cA{\dfrac{7}{20}}=\cK{0{,}35}$\par}\smallskip
{\centering\chip{fmB}{percentatge}\quad \chip{fmA}{fracció}\quad \chip{fmK}{decimal (tant per u)}\par}\smallskip
El $100\pc$ és el \textbf{total} ($=1$). Més del $100\pc$ és \textbf{més} que el total.
\end{idea}
\end{minipage}

\vfill
\subsection{Com passem d'una forma a una altra}
\begin{center}
\begin{tikzpicture}[b/.style={draw=#1,fill=#1!12,line width=1.3pt,rounded corners=3mm,minimum width=3.6cm,minimum height=1.3cm,align=center,font=\large}]
  \node[b=fmB] (p) at (0,0) {\heavy Percentatge\\$35\pc$};
  \node[b=fmK] (d) at (7,0) {\heavy Decimal\\$0{,}35$};
  \node[b=fmA] (f) at (3.5,-2.6) {\heavy Fracció\\$\dfrac{35}{100}=\dfrac{7}{20}$};
  \draw[-{Latex[length=3mm]},line width=1.3pt,fmmuted] ([yshift=5pt]p.east) -- node[above,font=\comic]{dividim entre 100} ([yshift=5pt]d.west);
  \draw[-{Latex[length=3mm]},line width=1.3pt,fmmuted] ([yshift=-5pt]d.west) -- node[below,font=\comic]{multipliquem per 100} ([yshift=-5pt]p.east);
  \draw[-{Latex[length=3mm]},line width=1.3pt,fmmuted] (p.south) -- node[left=4pt,font=\comic,align=right]{$\dfrac{p}{100}$ i\\simplifiquem} (f.west);
  \draw[-{Latex[length=3mm]},line width=1.3pt,fmmuted] (f.east) -- node[right=4pt,font=\comic,align=left]{dividim\\dalt entre baix} (d.south);
\end{tikzpicture}
\end{center}

\vfill
\begin{exemple}[Exemples]
\begin{minipage}[t]{0.48\linewidth}
$4\pc=0{,}04\qquad 120\pc=1{,}2$\\[3pt]
$0{,}375=37{,}5\pc\qquad 1{,}35=135\pc$
\end{minipage}\hfill
\begin{minipage}[t]{0.49\linewidth}
$60\pc=\dfrac{60}{100}=\dfrac35$\\[3pt]
$\dfrac38=0{,}375=37{,}5\pc\qquad \dfrac{7}{25}=\dfrac{28}{100}=28\pc$
\end{minipage}
\end{exemple}

\vfill
\begin{center}
\taularod{>{\columncolor{fmblue}\color{white}\heavy}l*{9}{w{c}{1.08cm}}}{
\rowcolor{fmblue}\% & \color{white}$1\pc$ & \color{white}$5\pc$ & \color{white}$10\pc$ & \color{white}$20\pc$ & \color{white}$25\pc$ & \color{white}$33{,}\overline{3}\pc$ & \color{white}$50\pc$ & \color{white}$75\pc$ & \color{white}$100\pc$\\
Fracció & $\frac{1}{100}$ & $\frac{1}{20}$ & $\frac{1}{10}$ & $\frac15$ & $\frac14$ & $\frac13$ & $\frac12$ & $\frac34$ & $1$\\
\rowcolor{fmsurface}\cellcolor{fmblue}Decimal & $0{,}01$ & $0{,}05$ & $0{,}1$ & $0{,}2$ & $0{,}25$ & $0{,}\overline{3}$ & $0{,}5$ & $0{,}75$ & $1$
}
\end{center}
{\comic\color{fmblue}\centering Val la pena saber-la de memòria: el $50\pc$ és la meitat, el $25\pc$ la quarta part, el $10\pc$ la desena part\dots\par}

\vfill
\clearpage
\subsection{Càlcul de percentatges}

\vfill
\begin{minipage}[c]{0.52\linewidth}
\begin{tcolorbox}[enhanced,colback=fmyellow!18,colframe=fmgold,arc=4mm,boxrule=1.6pt,before upper=\centering,top=3mm,bottom=3mm]
{\Large$p\pc$ de $C=C\cdot\dfrac{p}{100}$}\par\medskip
$15\pc$ de $80=80\cdot0{,}15=\boxed{12}$
\end{tcolorbox}
\end{minipage}\hfill
\begin{minipage}[c]{0.45\linewidth}
\begin{idea}[És una regla de tres]
{\centering$\left.\begin{array}{l}100\ \longrightarrow\ p\\ C\ \ \ \longrightarrow\ x\end{array}\right\}\quad x=\dfrac{C\cdot p}{100}$\par}
\end{idea}
\end{minipage}

\vfill
\subsection{Trucs de càlcul mental}
\begin{tcbraster}[raster columns=4,raster equal height,raster column skip=8pt]
\begin{tcolorbox}[fmbase,colback=fmdefbg,colframe=fmdef,before upper=\centering,top=2mm,bottom=2mm]
{\Large\heavy\color{fmblue}$50\pc$}\par la \textbf{meitat}\par{\small$50\pc$ de $84=42$}
\end{tcolorbox}
\begin{tcolorbox}[fmbase,colback=fmdefbg,colframe=fmdef,before upper=\centering,top=2mm,bottom=2mm]
{\Large\heavy\color{fmblue}$25\pc$}\par la meitat de la meitat\par{\small$25\pc$ de $84=21$}
\end{tcolorbox}
\begin{tcolorbox}[fmbase,colback=fmdefbg,colframe=fmdef,before upper=\centering,top=2mm,bottom=2mm]
{\Large\heavy\color{fmblue}$10\pc$}\par dividim entre 10\par{\small$10\pc$ de $84=8{,}4$}
\end{tcolorbox}
\begin{tcolorbox}[fmbase,colback=fmdefbg,colframe=fmdef,before upper=\centering,top=2mm,bottom=2mm]
{\Large\heavy\color{fmblue}$1\pc$}\par dividim entre 100\par{\small$1\pc$ de $84=0{,}84$}
\end{tcolorbox}
\end{tcbraster}

\vfill
\begin{minipage}[t]{0.5\linewidth}
\begin{truc}[Combina'ls]
$5\pc$: la meitat del $10\pc$. \ $20\pc$: el doble del $10\pc$.\\[4pt]
$15\pc$ de $60=\underbrace{6}_{10\pc}+\underbrace{3}_{5\pc}=9$
\end{truc}
\end{minipage}\hfill
\begin{minipage}[t]{0.47\linewidth}
\begin{truc}[Es poden intercanviar]
El $p\pc$ de $C$ és igual que el $C\pc$ de $p$:\\[4pt]
el $8\pc$ de $50$ = el $50\pc$ de $8=\boxed{4}$
\end{truc}
\end{minipage}

\vfill
\begin{exemple}[Exemple: \ en un institut de 640 alumnes, el $45\pc$ són xiques. Quantes n'hi ha?]
$45\pc$ de $640=640\cdot0{,}45=\boxed{288}$ xiques \qquad \nota{o bé: $\ 10\pc=64,\ \ 40\pc=256,\ \ 5\pc=32\ \to\ 256+32=288$}
\end{exemple}

\vfill
% ============================================================ 5.10
\clearpage
\seccio{sec10}{Augments i disminucions percentuals}{Multipliquem directament per l'índex de variació.}

\vfill
\begin{center}
\begin{tikzpicture}[x=0.075cm]
  \filldraw[fill=fmA!25,draw=fmA,line width=1.2pt] (0,0) rectangle (100,0.8); \node at (50,0.4) {\large$100\pc$ \ (inicial)};
  \filldraw[fill=fmK!30,draw=fmK,line width=1.2pt] (100,0) rectangle (115,0.8); \node at (107.5,0.4) {$+15$};
  \draw[decorate,decoration={brace,amplitude=5pt},fmK,line width=1.1pt] (0,0.95) -- (115,0.95) node[midway,above=6pt,font=\large,text=fmK!80!black]{augment del $15\pc$: \ $115\pc\ \to\ I_V=1{,}15$};
  \begin{scope}[yshift=-1.9cm]
  \filldraw[fill=fmA!25,draw=fmA,line width=1.2pt] (0,0) rectangle (80,0.8); \node at (40,0.4) {\large$80\pc$ \ (queda)};
  \filldraw[fill=fmR!8,draw=fmR,line width=1.2pt,dashed] (80,0) rectangle (100,0.8); \node[text=fmR] at (90,0.4) {$-20$};
  \draw[decorate,decoration={brace,amplitude=5pt,mirror},fmR,line width=1.1pt] (0,-0.15) -- (80,-0.15) node[midway,below=6pt,font=\large,text=fmR]{disminució del $20\pc$: \ $80\pc\ \to\ I_V=0{,}8$};
  \end{scope}
\end{tikzpicture}
\end{center}

\vfill
\begin{tcolorbox}[enhanced,colback=fmyellow!18,colframe=fmgold,arc=4mm,boxrule=1.6pt,before upper=\centering,top=3mm,bottom=3mm]
{\Large$C_f=C_i\cdot I_V$}\qquad\qquad augment: \ $I_V=1+\dfrac{p}{100}$\qquad\quad disminució: \ $I_V=1-\dfrac{p}{100}$
\end{tcolorbox}

\vfill
\begin{minipage}[t]{0.5\linewidth}
\begin{exemple}[Exemples]
Un preu de 80 € puja un $15\pc$:\\[2pt]
\hspace*{1em}$80\cdot1{,}15=\boxed{92\text{ €}}$\\[4pt]
Una jaqueta de 50 € té un $20\pc$ de descompte:\\[2pt]
\hspace*{1em}$50\cdot0{,}8=\boxed{40\text{ €}}$
\end{exemple}
\end{minipage}\hfill
\begin{minipage}[t]{0.47\linewidth}
\begin{idea}[Llegir un índex de variació]
$I_V>1$: \textbf{augmenta}. \ $I_V<1$: \textbf{disminueix}.\\[4pt]
$1{,}15\to\cK{+15\pc}\qquad 0{,}92\to\cR{-8\pc}$\\[3pt]
$1{,}05\to\cK{+5\pc}\qquad\ \ 0{,}7\to\cR{-30\pc}$\\[3pt]
$2\to\cK{+100\pc}$
\end{idea}
\end{minipage}

\vfill
\begin{minipage}[t]{0.5\linewidth}
\begin{truc}[I si coneixem l'inici i el final?]
$I_V=\dfrac{C_f}{C_i}$: \ de 40 € a 46 €, \ $\dfrac{46}{40}=1{,}15\to+15\pc$
\end{truc}
\end{minipage}\hfill
\begin{minipage}[t]{0.47\linewidth}
\begin{compte}[Un sol pas]
No cal calcular l'augment i sumar-lo: multiplica per $I_V$ i ja ho tens.
\end{compte}
\end{minipage}

\vfill
\begin{exemple}[Més exemples]
\begin{minipage}[t]{0.48\linewidth}
\textbf{IVA del $21\pc$}: un article val 50 € sense IVA.\\[3pt]
$50\cdot1{,}21=\boxed{60{,}50\text{ €}}$
\end{minipage}\hfill
\begin{minipage}[t]{0.49\linewidth}
Un poble de 25\,000 habitants perd el $4\pc$.\\[3pt]
$25\,000\cdot0{,}96=\boxed{24\,000\text{ habitants}}$
\end{minipage}
\end{exemple}

\vfill
% ============================================================ resum
\resumtitol
\begin{tcbraster}[raster columns=2,raster equal height=rows,raster column skip=14pt,raster row skip=20pt,raster before skip=12pt]
\rtile{5.1}{Decimals}{Part entera i part decimal.\\[1ex] Dècimes, centèsimes, mil·lèsimes.}
\rtile{5.2}{Comparar i aproximar}{Igualem el nombre de decimals.\\[1ex] Arrodonir: 5 o més, sumem 1.}
\rtile{5.3}{Sumar i multiplicar}{Sumar: comes alineades.\\[1ex] Producte: sumem els decimals.}
\rtile{5.4}{Dividir}{Divisor decimal: $\times10$, $\times100$\dots\\[1ex] fins que siga natural.}
\rtile{5.5}{Decimals i fraccions}{Exacte, periòdic pur o periòdic mixt.}
\rtile{5.6}{Proporcions}{$\frac ab=\frac cd\iff a\cdot d=b\cdot c$.}
\rtile{5.7}{Directa}{El doble, el doble.\\[1ex] Es manté el quocient.}
\rtile{5.8}{Inversa}{El doble, la meitat.\\[1ex] Es manté el producte.}
\rtile{5.9}{Percentatges}{$x\,\%$ de $C$: $C\cdot\frac{x}{100}$.}
\rtile{5.10}{Augments i descomptes}{Multipliquem per $1\pm\frac{x}{100}$.}
\end{tcbraster}
\end{document}
